\chapter{Analyse de besoins et méthodologie}

\section{I. Spécification des besoins et des données}

\subsection{Analyse et spécification des besoins}

Les utilisateurs finaux visés par ce projet sont les particuliers souhaitant estimer la valeur
d'un bien avant une transaction, ainsi que les professionnels (agences immobilières) souhaitant
disposer d'un outil d'aide à la décision rapide.

\textbf{Besoins fonctionnels :}
\begin{itemize}
  \item Estimer le prix d'un bien résidentiel (maison ou appartement) à partir de ses
  caractéristiques principales (type, quartier, surface, nombre de chambres).
  \item Comparer les résultats de plusieurs modèles pour un même bien.
  \item Explorer visuellement le jeu de données ayant servi à l'entraînement (distribution des
  prix, prix par quartier, etc.).
  \item Visualiser les performances comparées des modèles.
\end{itemize}

\textbf{Besoins non fonctionnels :}
\begin{itemize}
  \item Accessibilité via une interface web, sans compétence technique requise.
  \item Reproductibilité du pipeline de collecte, nettoyage et entraînement.
  \item Transparence sur la marge d'erreur des estimations produites.
\end{itemize}

\subsection{Sources de données}

Une recherche préalable a permis d'écarter plusieurs pistes de données existantes :
\begin{itemize}
  \item aucun jeu de données immobilier sénégalais n'est disponible sur les plateformes usuelles
  (Kaggle) ;
  \item l'Agence Nationale de la Statistique et de la Démographie (ANSD) publie des indices
  macro-économiques du secteur (indice du coût de la construction, indice des prix des
  matériaux), mais aucune donnée au niveau transaction ;
  \item aucun portail cadastral ouvert, comparable à celui existant en France, n'est disponible
  pour le Sénégal.
\end{itemize}

Faute de source structurée existante, les données ont été collectées par extraction
automatisée (\emph{web scraping}) depuis \textbf{Expat-Dakar}
(\url{https://www.expat-dakar.com/}), plateforme d'annonces généralistes hébergeant plusieurs
milliers d'annonces immobilières actives. Ce choix a été validé par une double vérification :
la structure des pages (attributs \texttt{data-t-listing\_*} intégrés directement dans le HTML,
encodant prix, type de bien, localisation et identifiant de façon structurée) permet une
extraction fiable, et le fichier \texttt{robots.txt} du site autorise l'accès aux pages
d'annonces (seuls les répertoires \texttt{/admin/}, \texttt{/compte/} et \texttt{/cdn-cgi/} sont
exclus).

\subsection{Processus de collecte}

La collecte a été réalisée avec un script Python (\texttt{src/scraper.py}) s'appuyant sur les
bibliothèques \texttt{requests} et \texttt{BeautifulSoup}. Le périmètre retenu couvre :
\begin{itemize}
  \item les annonces de \textbf{vente} (à l'exclusion de la location, pour ne pas mélanger deux
  marchés aux dynamiques de prix différentes) ;
  \item les catégories \textbf{maisons à vendre} et \textbf{appartements à vendre} (à
  l'exclusion des terrains nus et des biens commerciaux, hors périmètre résidentiel) ;
  \item l'ensemble du territoire sénégalais (pas de restriction géographique a priori).
\end{itemize}

Pour chaque annonce, les champs suivants sont extraits : identifiant, titre, catégorie, prix
(FCFA), quartier, ville/région, nombre de chambres, surface (m\textsuperscript{2}), date de
publication et URL. La pagination est parcourue jusqu'à réception d'une page vide ou d'une
réponse HTTP 404 signalant la fin des résultats, avec un délai d'une seconde et demie entre
chaque requête afin de limiter la charge sur le serveur cible. La collecte a produit
\textbf{660 annonces brutes} (330 maisons, 330 appartements), sans doublon.

\subsection{Pré-processing des données}

Le nettoyage du jeu de données brut (détaillé et exécuté dans le notebook
\texttt{01\_eda\_nettoyage.ipynb}) comprend les étapes suivantes :

\begin{enumerate}
  \item \textbf{Suppression des prix non exploitables.} Un seuil minimal de 5\,000\,000 FCFA a
  été retenu après inspection manuelle des valeurs basses : en deçà, les annonces correspondent
  systématiquement à des incohérences flagrantes plutôt qu'à de véritables biens bon marché : par
  exemple, une villa de 1000 m\textsuperscript{2} au quartier des Ngor (où le prix moyen dépasse
  200 millions FCFA) affichée à 1,5 million FCFA, probablement un prix « sur demande » ou une
  mensualité mal étiquetée comme prix total. Cette étape a supprimé 43 annonces.
  \item \textbf{Exclusion des biens hors périmètre résidentiel unitaire.} Les annonces présentant
  une surface supérieure à 2000 m\textsuperscript{2} ou plus de 15 chambres ont été exclues, ces
  valeurs correspondant vraisemblablement à des immeubles entiers plutôt qu'à un logement
  individuel. Cette étape a supprimé 14 annonces supplémentaires.
  \item \textbf{Imputation des valeurs manquantes.} Les champs \texttt{surface\_m2} et
  \texttt{bedrooms}, absents respectivement pour 16{,}4~\% et 23{,}3~\% des annonces (non
  systématiquement renseignés par les vendeurs), ont été imputés par la médiane du groupe
  (catégorie $\times$ ville/région), avec repli sur la médiane globale en l'absence de groupe
  suffisant.
  \item \textbf{Regroupement des quartiers rares.} Le jeu de données comporte 92 quartiers
  distincts pour 603 lignes, avec une forte hétérogénéité de représentation. Les quartiers
  comptant moins de 5 annonces ont été regroupés dans une catégorie \texttt{Autre} afin de
  limiter le risque de surapprentissage lors de l'encodage catégoriel, ramenant la cardinalité
  effective à 34 catégories.
\end{enumerate}

Le jeu de données final compte \textbf{603 annonces}, décrites par : le type de bien
(catégorie), le quartier (regroupé), la surface et le nombre de chambres, pour une cible à
prédire (le prix en FCFA).

\section{II. Méthodologie}

\subsection{Choix des algorithmes et modèles}

Quatre modèles ont été retenus et comparés :
\begin{enumerate}
  \item \textbf{Régression linéaire} : modèle de référence, non optimisé, permettant de juger
  de l'apport réel des modèles plus complexes.
  \item \textbf{Random Forest} \citep{breiman2001randomforests} : méthode ensembliste à base
  d'arbres, optimisée.
  \item \textbf{XGBoost} \citep{chen2016xgboost} : méthode de \emph{gradient boosting},
  optimisée.
  \item \textbf{MLP (Multi-Layer Perceptron)} : réseau de neurones artificiel, optimisé,
  représentant l'approche deep learning du projet.
\end{enumerate}

\subsection{Justification des choix}

La régression linéaire sert de ligne de base simple et interprétable, dans la tradition du
modèle hédonique. Random Forest et XGBoost ont été retenus car la littérature
(\S~2.2) montre leur supériorité constante sur les modèles linéaires pour ce type de problème,
tout en restant robustes sur des jeux de données de taille modeste. Le MLP a été retenu pour
répondre explicitement à l'objectif de comparaison avec une approche de deep learning formulé
dans le sujet du projet.

Un choix d'implémentation mérite d'être justifié : le réseau de neurones a été implémenté avec
la classe \texttt{MLPRegressor} de scikit-learn plutôt qu'avec un framework de deep learning
dédié (TensorFlow/Keras, utilisé par exemple dans des projets de prévision météorologique sur
des séries beaucoup plus longues). Avec environ 600 lignes d'entraînement, un framework complet
(gestion GPU, callbacks avancés) n'apporte aucun bénéfice ; un perceptron multicouche reste le
niveau de complexité approprié, tout en permettant une véritable optimisation d'hyperparamètres
(architecture, régularisation, taux d'apprentissage).

\subsection{Paramètres et configuration}

Chaque modèle (à l'exception de la régression linéaire) a été optimisé par
\texttt{RandomizedSearchCV} (recherche aléatoire, 30 combinaisons tirées, validation croisée à
5 plis mélangés), avec comme métrique d'optimisation l'erreur absolue moyenne négative
(\texttt{neg\_mean\_absolute\_error}) calculée sur la cible transformée en logarithme (voir
\S~3.5). Les espaces de recherche retenus sont résumés dans le tableau~\ref{tab:hyperparams}.

\begin{table}[ht]
\centering
\caption{Espaces d'hyperparamètres explorés par modèle}
\label{tab:hyperparams}
\small
\begin{tabular}{@{}ll>{\raggedright\arraybackslash}p{6.5cm}@{}}
\toprule
\textbf{Modèle} & \textbf{Hyperparamètre} & \textbf{Valeurs explorées} \\
\midrule
\multirow{5}{*}{Random Forest}
 & n\_estimators        & 100, 200, 300, 400, 500 \\
 & max\_depth            & None, 5, 10, 15, 20, 30 \\
 & min\_samples\_split   & 2, 5, 10 \\
 & min\_samples\_leaf    & 1, 2, 4 \\
 & max\_features         & sqrt, log2, None \\
\midrule
\multirow{7}{*}{XGBoost}
 & n\_estimators        & 100, 200, 300, 400 \\
 & max\_depth            & 3, 4, 5, 6, 8 \\
 & learning\_rate        & 0.01, 0.03, 0.05, 0.1, 0.2 \\
 & subsample             & 0.6, 0.8, 1.0 \\
 & colsample\_bytree     & 0.6, 0.8, 1.0 \\
 & reg\_alpha            & 0, 0.1, 1 \\
 & reg\_lambda           & 1, 1.5, 2 \\
\midrule
\multirow{4}{*}{MLP}
 & hidden\_layer\_sizes  & (32,), (64,), (64,32), (128,64), (64,64,32) \\
 & alpha                 & 0.0001, 0.001, 0.01, 0.1 \\
 & learning\_rate\_init  & 0.0005, 0.001, 0.005, 0.01 \\
 & activation            & relu, tanh \\
\bottomrule
\end{tabular}
\end{table}

\subsection{Évaluation des performances}

L'évaluation est réalisée sur un jeu de test indépendant, obtenu par partitionnement
aléatoire (80~\% entraînement / 20~\% test, soit 482 et 121 annonces respectivement), non
utilisé lors de l'optimisation des hyperparamètres.

\subsection{Métriques utilisées}

Quatre métriques complémentaires sont calculées, après retransformation des prédictions de
l'échelle logarithmique vers le FCFA :

\begin{itemize}
  \item \textbf{RMSE} (\emph{Root Mean Squared Error}) : $\sqrt{\frac{1}{n}\sum_i (y_i -
  \hat{y}_i)^2}$, pénalise fortement les grandes erreurs.
  \item \textbf{MAE} (\emph{Mean Absolute Error}) : $\frac{1}{n}\sum_i |y_i - \hat{y}_i|$,
  erreur moyenne absolue, en FCFA.
  \item \textbf{MAPE} (\emph{Mean Absolute Percentage Error}) : $\frac{100}{n}\sum_i
  \left|\frac{y_i - \hat{y}_i}{y_i}\right|$, erreur moyenne relative, la plus interprétable pour
  un public non technique.
  \item \textbf{$R^2$} : part de variance de la cible expliquée par le modèle.
\end{itemize}

\subsection{Protocole d'évaluation}

La distribution des prix étant fortement asymétrique (quelques biens haut de gamme, cf.
Figure~\ref{fig:eda-distribution}), tous les modèles sont entraînés sur la cible transformée en
$\log(1+prix)$, une pratique usuelle en modélisation de prix immobiliers qui stabilise la
variance et améliore la performance des modèles linéaires et des réseaux de neurones, sans
pénaliser les méthodes à base d'arbres, insensibles aux transformations monotones. Les
prédictions sont retransformées en FCFA (fonction exponentielle inverse) avant le calcul des
métriques finales, afin de conserver une interprétation directe des résultats.

\section{Conclusion}

Ce chapitre a détaillé la démarche méthodologique retenue : faute de données existantes, un
jeu de données a été constitué par extraction automatisée puis nettoyé selon des règles
justifiées et quantifiées ; quatre modèles ont été sélectionnés pour couvrir à la fois
l'approche hédonique classique, les méthodes ensemblistes et le deep learning, avec un
protocole d'optimisation et d'évaluation homogène permettant une comparaison rigoureuse,
présentée au chapitre suivant.
